\documentclass[10pt,a4paper]{amsart}
\usepackage[margin=19mm]{geometry}
\usepackage{amsmath,amssymb,mathtools}
\usepackage[scr=rsfs,cal=euler]{mathalfa}
\usepackage{xfrac}
\usepackage[unicode,hidelinks,bookmarks=false]{hyperref}
\newcommand{\ie}{i.e.\ }
\newcommand{\cf}{cf.\ }
\newcommand{\resp}{respectively\ }
\newcommand{\Subsection}{\subsection}
\providecommand{\nobreakdash}{\nobreak\hbox{-}}
\setlength{\emergencystretch}{2em}
\multlinegap0pt
\usepackage{amssymb}
\usepackage{mathtools}
\newtheorem{lemma}{Lemma}
\title{Closed Formulas for $\eta$-Corrections in the Once Punctured Torus}
\author{Nelson A. Colon Vargas}
\date{Source: arXiv:2508.18334v3 (CC BY 4.0)}
\begin{document}
\maketitle

\noindent\emph{Source and licence.} Closed Formulas for $\eta$-Corrections in the Once Punctured Torus. Version arXiv:2508.18334v3 (CC BY 4.0), distributed by arXiv under the Creative Commons Attribution 4.0 International licence, http://creativecommons.org/licenses/by/4.0/, as recorded in the arXiv licence metadata for this exact version and shown on the abstract page https://arxiv.org/abs/2508.18334v3. Full licence notice: \texttt{licenses/arxiv-2508.18334-CC-BY-4.0.txt}.\par\smallskip

\noindent\textit{Editorial scope.} Counted unit: the article's lemma lem:cascade (source0.tex lines 657--668). The article's complete original proof of it is delivered with it (source0.tex lines 670--685, one \texttt{proof} environment of 16 lines). No mathematics is written, repaired or re-derived: the statement and the proof are the article's own lines, in the article's own notation, and no retained line is substituted or re-flowed.  No further source-local context is needed: every cross-reference of every retained line resolves inside the two delivered spans.  Nothing was omitted from any retained span, so the package carries no omission record.  No retained line contains a citation, so the wrapper carries no bibliography and no bibliography entry.  No retained line contains a figure environment, an image-inclusion command or drawing code, so no image is delivered and the package declares no asset dependency.  Editorial formatting only, in addition: an AMS \texttt{amsart} preamble that restates the article's own declarations and macros used by the retained lines, namely the environments \texttt{lemma} on separate counters, together with the standard packages those lines need.  The retained environments print in this document as Lemma 1 in order of appearance, whereas the article prints the same items with the numbering of its own full document.  The complete native source of the article as distributed in the arXiv e-print for this exact version is bundled unchanged as \texttt{sources/183858/source0.tex}.
\begin{lemma}[Chebyshev cascade identity]\label{lem:cascade}
Let $\mu=(0,1)$ and put $\varepsilon=\operatorname{sgn}(n)$, $x:=t^2+t^{-2}$.
Define
\[
G_n \;:=\;\sum_{j=0}^{\lfloor (n-2)/2\rfloor}
t^{\varepsilon(n-2-2j)}\Big(T_{n-2-2j}(\mu)-\delta_{n-2-2j,0}\Big)\,S_j(x).
\]
Then
\begin{equation}\label{eq:cascade}
\mu_T\ast G_n \;=\; t^{\varepsilon}\,G_{n-1} \;+\; t^{-\varepsilon}\,G_{n+1} \;-\; t^{-\varepsilon n}\,.
\end{equation}
\end{lemma}

\begin{proof}
Write $k=n-2-2j$ (so $k\ge0$ and $k\equiv n\bmod 2$). Using
\[
\mu_T T_k(\mu)=T_{k+1}(\mu)+T_{k-1}(\mu)\quad(k\ge1),\qquad
S_{j+1}(x)=xS_j(x)-S_{j-1}(x),\quad S_0=1,\ S_1=x,
\]
we have, termwise,
\[
\mu_T\!\left(t^{\varepsilon k}\,[T_k(\mu)-\delta_{k,0}]\,S_j(x)\right)
= t^{\varepsilon k}\big(T_{k+1}+T_{k-1}-\delta_{k,0}\big)S_j(x).
\]
Reindex the $j$-sum: the $T_{k+1}$ part is the $j$th summand of $t^{\varepsilon}G_{n-1}$; the $T_{k-1}$ part
is the $(j+1)$st summand of $t^{-\varepsilon}G_{n+1}$. The only residue is when $k=0$,
where $T_0(\mu)=2$ and we have subtracted $1$; this yields the unit $1=T'_0$ and produces the final
$-\,t^{-\varepsilon n}$ in \eqref{eq:cascade}. Summing over $j$ gives the claim.
\end{proof}

\end{document}
